\documentclass[11pt,a4paper]{article}

\usepackage[a4paper,margin=1.65cm,top=1.45cm,bottom=1.55cm]{geometry}
\usepackage{fontspec}
\usepackage[french]{babel}
\usepackage{microtype}
\usepackage{array}
\usepackage{tabularx}
\usepackage{booktabs}
\usepackage{ragged2e}
\usepackage[table]{xcolor}
\usepackage{enumitem}
\usepackage{titlesec}
\usepackage{fancyhdr}
\usepackage[most]{tcolorbox}
\usepackage{hyperref}
\usepackage{xurl}

\setmainfont{Arial}
\setsansfont{Arial}
\renewcommand{\familydefault}{\sfdefault}

\definecolor{primary}{HTML}{17365D}
\definecolor{secondary}{HTML}{2F75B5}
\definecolor{softblue}{HTML}{EAF3F8}
\definecolor{softgray}{HTML}{F3F5F7}
\definecolor{linegray}{HTML}{D5DCE3}
\definecolor{textgray}{HTML}{4C5967}
\definecolor{softgreen}{HTML}{EAF5EE}

\hypersetup{
  colorlinks=true,
  linkcolor=primary,
  urlcolor=secondary,
  pdftitle={Stack Android simple - logiciel de caisse},
  pdfauthor={AMH Solution}
}
\urlstyle{same}

\setlength{\parindent}{0pt}
\setlength{\parskip}{4pt}
\renewcommand{\arraystretch}{1.22}
\setlist[itemize]{leftmargin=5mm,itemsep=1.5pt,topsep=2pt}
\setlist[enumerate]{leftmargin=6mm,itemsep=1.5pt,topsep=2pt}

\titleformat{\section}
  {\Large\bfseries\color{primary}}
  {\thesection.}{0.45em}{}
\titlespacing*{\section}{0pt}{10pt}{4pt}

\titleformat{\subsection}
  {\normalsize\bfseries\color{secondary}}
  {}{0pt}{}
\titlespacing*{\subsection}{0pt}{7pt}{3pt}

\pagestyle{fancy}
\setlength{\headheight}{22pt}
\fancyhf{}
\fancyhead[L]{\small\color{textgray}AMH Solution --- logiciel de caisse}
\fancyhead[R]{\small\color{textgray}Stack Android simple}
\fancyfoot[C]{\small\color{textgray}\thepage}
\renewcommand{\headrulewidth}{0.35pt}
\renewcommand{\headrule}{\hbox to\headwidth{\color{linegray}\leaders\hrule height \headrulewidth\hfill}}

\newcolumntype{Y}{>{\RaggedRight\arraybackslash}X}
\newcolumntype{L}[1]{>{\RaggedRight\arraybackslash}p{#1}}
\newcommand{\source}[2]{\href{#2}{#1}}

\begin{document}

\begin{center}
  {\fontsize{22}{26}\selectfont\bfseries\color{primary}
  Stack Android simple pour APK}\\[3pt]
  {\large Version allégée pour logiciel de caisse restaurant}\\[7pt]
  {\small\color{textgray}Note technique courte --- 3 juin 2026}
\end{center}

\begin{tcolorbox}[
  colback=softgreen,
  colframe=primary,
  boxrule=0.8pt,
  arc=2pt,
  left=8pt,right=8pt,top=7pt,bottom=7pt
]
\textbf{\color{primary}Recommandation}\\
Pour le client actuel et un devis limité, partir sur une \textbf{application Android native en APK}
avec \textbf{Kotlin}, \textbf{Jetpack Compose}, \textbf{Room / SQLite} et une impression validée
sur le matériel réel. Cette stack est plus simple à livrer qu'une architecture Electron ou cloud,
et correspond mieux à une tablette ou un terminal POS Android.
\end{tcolorbox}

\section{Objectif}

Livrer une première version exploitable centrée sur l'encaissement local :

\begin{itemize}
  \item installer l'application sous forme de fichier \texttt{.apk} ;
  \item fonctionner sans connexion Internet ;
  \item stocker les articles, commandes, paiements et historiques localement ;
  \item imprimer une note client après validation du matériel ;
  \item produire des exports simples en \texttt{JSON} ou \texttt{CSV}.
\end{itemize}

Cette proposition sert au périmètre client court. L'option interne Electron / React peut rester
une piste d'apprentissage ou de produit futur.

\section{Stack proposée}

\rowcolors{2}{softgray}{white}
\begin{tabularx}{\textwidth}{L{3.55cm} L{3.45cm} Y}
  \toprule
  \rowcolor{primary}
  \color{white}\textbf{Besoin} &
  \color{white}\textbf{Choix} &
  \color{white}\textbf{Rôle} \\
  \midrule
  Plateforme cible &
  Android &
  Installer l'application sur tablette Android ou terminal POS Android. \\
  Format de livraison &
  \texttt{APK} &
  Fournir un fichier installable directement sur l'appareil du client. \\
  Langage principal &
  Kotlin &
  Développer une application Android native maintenable. \\
  Interface tactile &
  Jetpack Compose &
  Créer des écrans tactiles modernes avec peu de code boilerplate. \\
  Base locale &
  Room / SQLite &
  Enregistrer articles, commandes, lignes, paiements et historiques localement. \\
  Configuration simple &
  DataStore &
  Stocker les paramètres légers : nom du restaurant, préférences, imprimante. \\
  Tickets et sauvegardes &
  JSON &
  Figer le contenu d'un ticket et exporter une sauvegarde lisible. \\
  Exports tableur &
  CSV &
  Fournir un format simple à ouvrir pour le responsable ou le comptable. \\
  Impression &
  SDK fabricant ou ESC/POS &
  Imprimer sur imprimante intégrée, Bluetooth, USB ou réseau après test réel. \\
  Build &
  Android Studio / Gradle &
  Compiler et générer l'APK. \\
  Tests &
  JUnit &
  Tester les calculs de panier, TVA, paiement mixte et rendu monnaie. \\
  Versioning &
  Git &
  Suivre les changements et sécuriser les livraisons. \\
  \bottomrule
\end{tabularx}
\rowcolors{2}{}{}

\vspace{5pt}
\textbf{Règle métier importante :} les montants sont stockés en centimes sous forme d'entiers,
par exemple \texttt{1850} pour 18,50 euros.

\section{Pourquoi cette stack est plus simple}

\rowcolors{2}{softgray}{white}
\begin{tabularx}{\textwidth}{L{4.5cm} Y Y}
  \toprule
  \rowcolor{primary}
  \color{white}\textbf{Critère} &
  \color{white}\textbf{APK Android native} &
  \color{white}\textbf{Electron / React interne} \\
  \midrule
  Installation client &
  Un APK à installer sur l'appareil &
  Application desktop Windows à empaqueter \\
  Matériel cible &
  Tablette ou terminal Android POS &
  Terminal Windows ou PC tactile \\
  Nombre de briques &
  Kotlin, Compose, Room &
  Electron, React, TypeScript, Vite, SQLite \\
  Coût matériel possible &
  Plus bas avec tablette Android simple &
  Souvent plus élevé sur tablette ou POS Windows \\
  Logique offline &
  Naturelle avec base locale sur l'appareil &
  Possible aussi, mais architecture plus riche \\
  \bottomrule
\end{tabularx}
\rowcolors{2}{}{}

\section{Périmètre fonctionnel simple}

\begin{itemize}
  \item catalogue d'articles et catégories ;
  \item panier et type de commande : sur place, à emporter, livraison ;
  \item paiements : espèces, carte confirmée manuellement, titres-restaurant, autre ;
  \item paiement mixte et rendu monnaie ;
  \item génération d'une note client ;
  \item historique des encaissements ;
  \item annulation ou correction traçable ;
  \item statistiques journalières simples ;
  \item export \texttt{JSON} et \texttt{CSV}.
\end{itemize}

\section{Points à valider avant chiffrage}

\begin{itemize}
  \item Le client possède-t-il déjà une tablette Android ou un terminal POS Android ?
  \item L'imprimante est-elle intégrée au terminal ou externe ?
  \item L'impression doit-elle passer par Bluetooth, USB, réseau ou SDK fabricant ?
  \item Le TPE bancaire reste-t-il autonome avec confirmation manuelle dans l'application ?
  \item L'installation de l'APK hors Play Store est-elle acceptée sur le matériel retenu ?
\end{itemize}

\begin{tcolorbox}[
  colback=softblue,
  colframe=secondary,
  boxrule=0.55pt,
  arc=2pt,
  left=8pt,right=8pt,top=6pt,bottom=6pt
]
\textbf{Point critique}\\
Le vrai risque de cette stack n'est pas Kotlin : c'est l'impression. Le choix final du terminal
doit donc être validé avec un test d'impression réel avant de promettre une livraison stable.
\end{tcolorbox}

\section{Plan d'action court}

\begin{enumerate}
  \item Créer le projet Android Kotlin dans Android Studio.
  \item Construire les premiers écrans Compose : accueil, catalogue, panier.
  \item Ajouter Room / SQLite : articles, commandes, lignes et paiements.
  \item Implémenter les calculs : total TTC, paiements mixtes, rendu monnaie.
  \item Générer un ticket JSON figé à chaque encaissement.
  \item Ajouter historique, correction et annulation traçable.
  \item Ajouter exports JSON et CSV.
  \item Tester l'impression sur le terminal ou l'imprimante retenue.
  \item Générer l'APK et tester l'installation sur l'appareil client.
  \item Réaliser une recette courte avec quelques scénarios d'encaissement.
\end{enumerate}

\section{Sources officielles}

{\footnotesize
\begin{itemize}[itemsep=0pt,topsep=0pt,leftmargin=4mm]
  \item Android Kotlin et Jetpack Compose :
  \source{documentation Android}{https://developer.android.com/jetpack/compose/kotlin}
  \item Jetpack Compose :
  \source{documentation Android}{https://developer.android.com/compose}
  \item Room / SQLite :
  \source{documentation Android}{https://developer.android.com/codelabs/basic-android-kotlin-compose-persisting-data-room}
  \item DataStore :
  \source{documentation Android}{https://developer.android.com/datastore}
\end{itemize}
}

\vfill
{\small\color{textgray}
Document de travail. Cette stack vise une version plus simple et moins coûteuse pour un client
précis. Les choix d'impression, de terminal et d'installation doivent être validés sur matériel réel.
}

\end{document}
